\section{Statistical details}\label{app:tables}

\subsection*{Paired tests and multiplicity}

The two primary tests of Section~\ref{sec:system} are the DeLong comparison of the
linear model against the gradient-boosted gate and the exact McNemar comparison of
their paired decisions. DeLong's statistic gives $\Delta$AUC $=\DelongDelta$ with
standard error $0.0078$ and $z=\DelongZ$; McNemar counts $b=\McNemarB$ rows
correct for the gate and wrong for the linear model against $c=\McNemarC$ in the
opposite direction (\McNemarDiscordant{} discordant pairs). Holm's step-down
procedure over the two raw $p$-values rejects both at $\alpha=0.05$, and the
adjusted values equal the raw ones to the precision we store. We report the
adjusted values and the effect sizes rather than the raw $p$-values alone, because
at this sample size a $p$-value below $10^{-3}$ is uninformative about the
magnitude of the difference.

\subsection*{Bootstrap units}

Two resampling units are reported for every paired difference: the row and the
entity (card). With \BootstrapB{} resamples the row-level interval for the gate's
AUC-PR is \ClusterCiLo--\ClusterCiHi{} and the time-block interval
(\BootstrapBlocks{} blocks of 512 rows) is \BlockCiLo--\BlockCiHi. The entity-level
interval is wider than the row-level one in every pair we compute, which is the
expected direction: rows that share a card are not independent, so the row unit
understates the variance. We report both and draw conclusions from the entity-level
interval.

\subsection*{Power}

The number of training seeds is set by
$B=\big((z_{1-\alpha/2}+z_{\text{power}})\,\sigma/\delta\big)^2$ with
$\alpha=0.05$, power $0.8$, the observed seed standard deviation
$\sigma=\SeedSd$ and the observed seed-to-seed difference
$\delta=0.0069$. The result is \RequiredSeeds{} seeds, which is why the seed sweep
of \SeedN{} seeds is reported as a descriptive statistic (mean \SeedMean, standard
deviation \SeedSd) rather than as a hypothesis test. Reporting a Wilcoxon statistic
on five seeds, as the original pipeline did, would have produced a test with no
power to detect the effect it was testing for.

\subsection*{Challenge protocol}

The soundness of the commitment-and-challenge protocol is checked by Monte Carlo
over a grid of cheating fractions $f\in\{0.01,\ 0.05,\ 0.1,\ 0.2\}$ and
challenge rates $q\in\{0.01,\ 0.05,\ 0.1,\ 0.25\}$, \ChallengeCells{} cells with \ChalTrials{} trials each on
$N=\ChalN$ records. The observed escape frequencies follow the hypergeometric
bound closely: at $f=0.2$ and the smallest challenge rate the observed escape
frequency is \ChalEscapeObs{} against a theoretical \ChalEscapeTheory{}. The
detection rate rises from \ChalDetLow{} to \ChalDetHigh{} as the challenge
rate increases. The required challenge rate to force an escape probability below $0.01$
at the largest tolerable cheating fraction is $q=\ChallengeQ$. Replacing the
post-anchor unpredictable seed by a predictable one drives detection to zero, which
is the negative result that motivates the unpredictability requirement; the legacy
commitment scheme, which used a seeded nonce with low-entropy fields, was recovered
for every target it was applied to, which is why no zero-knowledge claim is made
anywhere in this paper.
